%=======================================================================
%  Foundations -- Design Documentation
%  CSE 360 Team Project
%
%  HOW TO USE THIS FILE
%  --------------------
%  1. Compile as-is. Every diagram shows a labelled placeholder box.
%  2. Draw each diagram in Astah using the specification printed next to
%     its placeholder.
%  3. Export each one as a PNG or PDF into a figures/ folder beside this
%     file, using the exact filename the placeholder names.
%  4. Recompile. The \diagram command detects the file and swaps it in
%     automatically -- you do not have to edit any LaTeX.
%
%  Only standard packages are used, so this compiles on Overleaf with no
%  configuration.
%=======================================================================

\documentclass[11pt,a4paper]{article}

\usepackage[T1]{fontenc}
\usepackage[utf8]{inputenc}
\usepackage[margin=1in]{geometry}
\usepackage{graphicx}
\usepackage{booktabs}
\usepackage{longtable}
\usepackage{tabularx}
\usepackage{array}
\usepackage{enumitem}
\usepackage{float}
\usepackage{xcolor}
\usepackage{caption}
\usepackage{listings}
\usepackage[hidelinks]{hyperref}

%-----------------------------------------------------------------------
% Colours
%-----------------------------------------------------------------------
\definecolor{accent}{HTML}{4338CA}
\definecolor{muted}{HTML}{667085}
\definecolor{codebg}{HTML}{F7F7FA}
\definecolor{rulecol}{HTML}{D0D5DD}

%-----------------------------------------------------------------------
% A diagram slot.
%   #1 caption   #2 filename in figures/   #3 label suffix
% Shows the exported file if it exists, otherwise a placeholder naming
% the file it is waiting for.
%-----------------------------------------------------------------------
\newcommand{\diagram}[3]{%
  \begin{figure}[H]
    \centering
    \IfFileExists{figures/#2}%
      {\includegraphics[width=\linewidth,height=0.42\textheight,keepaspectratio]{figures/#2}}%
      {\setlength{\fboxrule}{0.8pt}\color{rulecol}%
       \fbox{\parbox[c][0.26\textheight][c]{0.93\linewidth}{%
         \centering\color{muted}\itshape
         Diagram placeholder\\[0.7em]
         \normalfont\ttfamily\small figures/#2\\[0.7em]
         \itshape Draw this in Astah using the specification above,\\
         export it to that path, and recompile.}}}%
    \caption{#1}
    \label{fig:#3}
  \end{figure}}

%-----------------------------------------------------------------------
% "What to draw in Astah" heading
%-----------------------------------------------------------------------
\newcommand{\astah}[1]{%
  \par\medskip\noindent
  {\color{accent}\bfseries Astah --- #1}\par\nopagebreak\smallskip}

%-----------------------------------------------------------------------
% Code listing style
%-----------------------------------------------------------------------
\lstdefinestyle{java}{
  language=Java, basicstyle=\ttfamily\footnotesize,
  keywordstyle=\color{accent}, commentstyle=\color{muted}\itshape,
  stringstyle=\color{black}, backgroundcolor=\color{codebg},
  frame=single, rulecolor=\color{rulecol}, framesep=5pt,
  breaklines=true, showstringspaces=false, columns=fullflexible,
  xleftmargin=4pt, xrightmargin=4pt, aboveskip=10pt, belowskip=10pt}
\lstdefinestyle{sql}{
  language=SQL, basicstyle=\ttfamily\footnotesize,
  keywordstyle=\color{accent}, commentstyle=\color{muted}\itshape,
  backgroundcolor=\color{codebg}, frame=single,
  rulecolor=\color{rulecol}, framesep=5pt, breaklines=true,
  showstringspaces=false, columns=fullflexible,
  xleftmargin=4pt, xrightmargin=4pt}
\lstset{style=java}

\newcommand{\cls}[1]{\texttt{#1}}
\setlist[itemize]{topsep=4pt,itemsep=2pt,parsep=0pt}
\setlength{\parskip}{0.5em}
\setlength{\parindent}{0pt}

%=======================================================================
\begin{document}
%=======================================================================

\begin{titlepage}
\centering
\vspace*{4cm}
{\Huge\bfseries Foundations\par}
\vspace{0.6cm}
{\Large Design Documentation\par}
\vspace{1.2cm}
{\large CSE 360 --- Team Project\par}
\vspace{0.4cm}
{\color{muted}\large Architecture, detailed design, and traceability\par}
\vfill
{\color{muted}\small
This document describes the system as implemented.\\
Where the implementation is weak or incorrect, that is recorded in
Section~\ref{sec:issues} rather than corrected in the diagrams.\par}
\vspace{1cm}
\end{titlepage}

\tableofcontents
\newpage

%=======================================================================
\section{Architecture}
\label{sec:arch}
%=======================================================================

\subsection{Overview}

Foundations is a single-process JavaFX desktop application over an embedded H2 database.
There is no server, no network layer, and no concurrency: one person uses one window, and
the entire application shares a single database connection.

The system is organised in five layers.

\begin{center}
\begin{tabularx}{\linewidth}{@{}lX@{}}
\toprule
\textbf{Layer} & \textbf{Responsibility} \\
\midrule
Entry point & Connects to the database and chooses the first page \\
Presentation & Thirteen MVC page packages, one per screen \\
Shared UI services & Theme and layout builders; the password validation facade \\
Validation & Two finite state machines for username and password \\
Domain & The \cls{User} entity \\
Persistence & Every SQL statement in the system \\
\bottomrule
\end{tabularx}
\end{center}

\astah{Package diagram}
Create a package diagram containing the following packages, grouped into the five layers
above using package nesting or a note per layer.

\begin{itemize}
  \item \textbf{Entry:} \cls{applicationMain}
  \item \textbf{Presentation:} \cls{guiUserLogin}, \cls{guiFirstAdmin}, \cls{guiNewAccount},
        \cls{guiResetPassword}, \cls{guiAdminHome}, \cls{guiListUsers},
        \cls{guiAddRemoveRoles}, \cls{guiDeleteUser}, \cls{guiOnetimePassword},
        \cls{guiMultipleRoleDispatch}, \cls{guiContributor}, \cls{guiViewer},
        \cls{guiCurator}, \cls{guiUserUpdate}
  \item \textbf{Shared UI:} \cls{guiTools}
  \item \textbf{Validation:} \cls{fPasswordPopUpWindow},
        \cls{userNameRecognizerTestbed}
  \item \textbf{Domain:} \cls{entityClasses}
  \item \textbf{Persistence:} \cls{database}
  \item \textbf{Test:} \cls{listUsersTesting}, \cls{fPasswordEvaluationTestbedMain}
\end{itemize}

Dependencies to draw, all as \textit{uses} arrows: every presentation package depends on
\cls{entityClasses}, \cls{database} and \cls{guiTools}; \cls{guiTools} depends on
\cls{fPasswordPopUpWindow}; \cls{guiFirstAdmin} depends on
\cls{userNameRecognizerTestbed}; \cls{database} depends on \cls{entityClasses};
\cls{applicationMain} depends on \cls{database}, \cls{guiFirstAdmin} and
\cls{guiUserLogin}.

\diagram{System architecture, by package and layer}{architecture-packages.png}{arch-pkg}

\subsection{Layer responsibilities}

\textbf{Entry point.} \cls{FoundationsMain} is the JavaFX \cls{Application}. It connects to
the database, then chooses the first page: if the database is empty, the person running the
application must be setting the system up, so they go to First Admin; otherwise they go to
Login. It also holds the only two pieces of genuinely global state --- the single
\cls{Database} instance and \cls{activeHomePage}, which records which role's home page to
return to.

\textbf{Presentation.} Thirteen packages, each a Model--View--Controller triad for one page.

\begin{itemize}
  \item The \textbf{View} is a singleton holding every widget as a \cls{static} field.
        \cls{displayXxx(...)} is the only entry point from outside the package; it builds
        the page on first use and refreshes per-visit values on every call.
  \item The \textbf{Controller} is a set of \cls{protected static} methods that the View's
        event handlers call. It reads values directly off the View's fields.
  \item The \textbf{Model} is an empty placeholder in most packages. Two are real:
        \cls{ModelListUsers} holds the user table's name and role formatting, and
        \cls{fPasswordPopUpWindow.Model} holds the password state machine.
\end{itemize}

\textbf{Shared UI services.} \cls{guiTools.Theme} loads the stylesheet and supplies the card,
header and footer builders every page is assembled from. \cls{guiTools.PasswordCheck} is the
single place the application asks whether a password is acceptable.

\textbf{Validation.} Two finite state machines, each translated from a directed-graph
diagram and each exercised by its own console testbed. See Section~\ref{sec:fsm}.

\textbf{Domain.} \cls{entityClasses.User} is the only domain type: one account, its name
fields, and its four role flags.

\textbf{Persistence.} \cls{database.Database} owns the connection and every SQL statement in
the system. No other class contains SQL.

\subsection{Navigation}

There is no router or navigation controller. A page moves to another page by calling that
page's \cls{displayXxx} method directly, so the packages form a dense call graph rather than
a tree.

\astah{Second package or component diagram, navigation only}
Draw each page package as a node and each navigation call as a directed arrow. The edges
are: Login $\rightarrow$ Admin Home, Role Dispatch, Reset Password; Role Dispatch
$\rightarrow$ Admin Home, Contributor, Viewer, Curator; Admin Home $\rightarrow$ List Users,
Add/Remove Roles, Delete User, One-Time Password, User Update; and each of List Users,
Add/Remove Roles, Delete User, One-Time Password and User Update $\rightarrow$ Admin Home.
Every page also has a Logout edge back to Login.

\diagram{Navigation between pages}{architecture-navigation.png}{arch-nav}

The consequence of this design is that a page cannot be moved or renamed without touching
every page that navigates to it, and the question ``where can the user get to from here''
is only answerable by searching the source.

\subsection{Architectural decisions and their consequences}

\begin{center}
\begin{tabularx}{\linewidth}{@{}p{0.30\linewidth}X@{}}
\toprule
\textbf{Decision} & \textbf{Consequence} \\
\midrule
Singleton Views with \cls{static} widget fields &
Cheap repeat visits and low memory use, but anything set in the constructor is set once for
the life of the application. Per-user values must be refreshed in \cls{displayXxx}. Getting
this wrong is not hypothetical --- see issue~\ref{iss:stale}. \\
\addlinespace
One \cls{Database} class for all persistence &
Easy to find where a query lives; also makes \cls{Database} the largest and most-modified
file in the project, with roughly forty public methods. \\
\addlinespace
Controllers read View fields directly &
Removes parameter passing, but couples each Controller to the exact field names of its View,
so a View cannot be rewritten without checking its Controller. \\
\addlinespace
Validation in testbed packages &
Both recognizers can be exercised without starting JavaFX, and the live application reaches
them through the same code the testbeds use, so the two cannot disagree. \\
\addlinespace
Peer-to-peer navigation &
No central place to change a flow; each page hard-codes its destinations. \\
\bottomrule
\end{tabularx}
\end{center}

\subsection{Constraints}

\begin{center}
\begin{tabularx}{\linewidth}{@{}p{0.34\linewidth}X@{}}
\toprule
\textbf{Constraint} & \textbf{Consequence} \\
\midrule
Single user, single window & No concurrency control anywhere; the singletons are safe only
because of this \\
Fixed $800\times600$ window & Every page must fit that height; there is no page-level
scrolling \\
Embedded H2, one connection & Only one instance of the application may run at a time \\
\cls{CREATE TABLE IF NOT EXISTS} & The schema is never migrated --- see
issue~\ref{iss:migrate} \\
No build tool & Dependencies are Eclipse user libraries, so the build is not reproducible
from the repository alone \\
\bottomrule
\end{tabularx}
\end{center}

\newpage
%=======================================================================
\section{Detailed design: classes}
\label{sec:classes}
%=======================================================================

\subsection{Domain and persistence}

\astah{Class diagram --- domain and persistence}
Two classes with a dependency from \cls{Database} to \cls{User}. Give \cls{User} all eleven
attributes and its accessors; give \cls{Database} the attributes and the operations grouped
below. Mark every \cls{Database} operation \cls{public} and every \cls{User} attribute
\cls{private}.

\diagram{Domain and persistence classes}{class-domain-persistence.png}{cls-domain}

\textbf{\cls{entityClasses.User}} --- one account.

\begin{center}
\begin{tabularx}{\linewidth}{@{}llX@{}}
\toprule
\textbf{Attribute} & \textbf{Type} & \textbf{Notes} \\
\midrule
\cls{userName} & String & Unique; cannot be changed after creation \\
\cls{password} & String & Stored in clear text --- see issue~\ref{iss:plaintext} \\
\cls{firstName} & String & \\
\cls{middleName} & String & May be null or empty \\
\cls{lastName} & String & \\
\cls{preferredFirstName} & String & May be null or empty; displayed instead of
\cls{firstName} when set \\
\cls{emailAddress} & String & \\
\cls{adminRole} & boolean & \\
\cls{contributorRole} & boolean & \\
\cls{viewerRole} & boolean & \\
\cls{curatorRole} & boolean & \\
\bottomrule
\end{tabularx}
\end{center}

Operations: a constructor taking all eleven values, a getter per attribute, setters for the
name and email fields and the four roles, and \cls{getNumRoles()} returning how many roles
the account plays.

\textbf{\cls{database.Database}} --- all persistence. Operations grouped by purpose:

\begin{center}
\begin{tabularx}{\linewidth}{@{}p{0.24\linewidth}X@{}}
\toprule
\textbf{Group} & \textbf{Operations} \\
\midrule
Connection &
\cls{connectToDatabase()}, \cls{closeConnection()}, \cls{createTables()} \\
\addlinespace
Accounts &
\cls{register(User)}, \cls{doesUserExist(String)}, \cls{deleteUser(String)},
\cls{getAllUserAccounts()}, \cls{getUserList()}, \cls{getNumberOfUsers()},
\cls{isDatabaseEmpty()} \\
\addlinespace
Authentication &
\cls{loginAdmin(User)}, \cls{loginContributor(User)}, \cls{loginViewer(User)},
\cls{loginCurator(User)}, \cls{getUserAccountDetails(String)} \\
\addlinespace
Roles &
\cls{updateUserRole(String, String, String)}, \cls{getNumberOfRoles(User)},
\cls{getNumberOfAdmins()}, \cls{getFoundingAdminUsername()},
\cls{isFoundingAdmin(String)} \\
\addlinespace
Profile fields &
\cls{getFirstName}/\cls{updateFirstName} and the same pair for middle name, last name,
preferred first name and email address \\
\addlinespace
Passwords &
\cls{updatePassword(String, String)},
\cls{updatePassword(String, String, boolean)},
\cls{getCurrentOnetimePasswordFlag()} \\
\addlinespace
Invitations &
\cls{generateInvitationCode(String, String)}, \cls{getNumberOfInvitations()},
\cls{emailaddressHasBeenUsed(String)}, \cls{getRoleGivenAnInvitationCode(String)},
\cls{getEmailAddressUsingCode(String)}, \cls{removeInvitationAfterUse(String)} \\
\bottomrule
\end{tabularx}
\end{center}

\cls{Database} also caches the most recently fetched account in a set of \cls{current...}
fields, populated by \cls{getUserAccountDetails}. Pages read those rather than re-querying.

\subsection{The page pattern}

Every page package follows the same three-class shape, so one diagram documents all
thirteen.

\astah{Class diagram --- the MVC page pattern}
Draw the pattern once using \cls{guiListUsers} as the concrete example, and add a note that
the remaining twelve page packages follow the same shape. Include:

\begin{itemize}
  \item \cls{ViewListUsers} --- \cls{+displayListUsers(Stage, User)},
        \cls{-refreshTable(String)}, \cls{-showSelection(User)},
        \cls{-changeRole(String, boolean)}, \cls{-deleteSelected()}
  \item \cls{ControllerListUsers} --- \cls{\#performRoleChange(String, String, boolean)},
        \cls{\#confirmDelete(String)}, \cls{\#performDelete(String)},
        \cls{\#performReturn()}, \cls{\#performLogout()}, \cls{\#performQuit()}
  \item \cls{ModelListUsers} --- \cls{+formatFullName(User)}, \cls{+formatRoles(User)}
\end{itemize}

Relationships: View $\rightarrow$ Controller (calls), Controller $\rightarrow$ View (reads
fields, so draw it bidirectional), View $\rightarrow$ Model (calls), View $\rightarrow$
\cls{Theme} (calls), View and Controller $\rightarrow$ \cls{Database} (calls).

\diagram{The Model--View--Controller page pattern}{class-mvc-pattern.png}{cls-mvc}

\subsection{Shared services}

\astah{Class diagram --- shared services}
Three classes, no inheritance, all operations \cls{static}.

\begin{center}
\begin{tabularx}{\linewidth}{@{}p{0.26\linewidth}X@{}}
\toprule
\textbf{Class} & \textbf{Operations} \\
\midrule
\cls{guiTools.Theme} &
\cls{apply(Scene)}, \cls{pageRoot(Node)}, \cls{card(double,double)},
\cls{pageHeader(Label,Label,Node)}, \cls{footerBar(Node...)}, \cls{label(String,String)},
\cls{button(String,String)}, \cls{style(Node,String)}, \cls{hairline()}, \cls{spacer()},
\cls{row(double,Node...)}, \cls{chip(String,boolean)} \\
\addlinespace
\cls{guiTools.PasswordCheck} &
\cls{firstProblem(String)}, \cls{-readableList(List)} \\
\addlinespace
\cls{guiListUsers.ModelListUsers} &
\cls{formatFullName(User)}, \cls{formatRoles(User)} \\
\bottomrule
\end{tabularx}
\end{center}

Dependency to draw: \cls{PasswordCheck} $\rightarrow$ \cls{fPasswordPopUpWindow.Model}.
This is the relationship that keeps the live application and the password testbed judging
input by the same rules.

\diagram{Shared services}{class-shared-services.png}{cls-shared}

\newpage
%=======================================================================
\section{Detailed design: state machines}
\label{sec:fsm}
%=======================================================================

Both recognizers are finite state machines translated by hand from directed-graph diagrams.
Each has a console testbed that drives it with fixed cases, so the diagrams below are
testable claims rather than descriptions.

\subsection{Username recognizer}

\cls{userNameRecognizerTestbed.UserNameRecognizer.checkForValidUserName(String)} returns an
empty string when the username is acceptable and a specific error message otherwise.

\astah{State machine diagram --- username recognizer}
Three states. State~1 is the only final state.

\begin{center}
\begin{tabularx}{\linewidth}{@{}clX@{}}
\toprule
\textbf{State} & \textbf{Meaning} & \textbf{Transitions} \\
\midrule
0 & Start & \texttt{A--Z, a--z} $\rightarrow$ 1. Any other character halts. \\
\addlinespace
1 & Final. A valid name so far &
\texttt{A--Z, a--z, 0--9} $\rightarrow$ 1 \newline
\texttt{-\ \_\ .\ \&} $\rightarrow$ 2 \newline
Any other character halts. \\
\addlinespace
2 & Just consumed a separator &
\texttt{A--Z, a--z, 0--9} $\rightarrow$ 1. Any other character halts. \\
\bottomrule
\end{tabularx}
\end{center}

Every transition increments a character counter, and the machine halts as soon as that
counter exceeds 16. Show the counter as an action on the transitions, and the length limit
as a guard.

\diagram{Username recognizer state machine}{fsm-username.png}{fsm-user}

The verdict when the machine halts depends on the state it stopped in:

\begin{center}
\begin{tabularx}{\linewidth}{@{}clX@{}}
\toprule
\textbf{Halt state} & \textbf{Condition} & \textbf{Result} \\
\midrule
0 & --- & Rejected: a username must start with \texttt{A--Z} or \texttt{a--z} \\
1 & fewer than 4 characters & Rejected: at least 4 characters required \\
1 & more than 16 characters & Rejected: no more than 16 characters \\
1 & input not fully consumed & Rejected: only \texttt{A--Z a--z 0--9 - \_ . \&} allowed \\
1 & 4--16 characters, fully consumed & \textbf{Accepted} \\
2 & --- & Rejected: the character after a separator must be \texttt{A--Z a--z 0--9} \\
\bottomrule
\end{tabularx}
\end{center}

In plain terms: a username must start with a letter, be 4 to 16 characters long, contain
only letters, digits and the separators \texttt{-\ \_\ .\ \&}, and never end on a separator
or place two separators together.

\subsection{Password recognizer}

\cls{fPasswordPopUpWindow.Model.evaluatePassword(String)} walks the input once, setting a
flag for each requirement it satisfies, and halts immediately on a character it does not
recognise.

\astah{State machine diagram --- password recognizer}
A single scanning state with a self-transition per accepted character class, plus an error
exit. Model it as: an initial state; a \textit{Scanning} state; self-transitions labelled
\texttt{A--Z / set foundUpperCase}, \texttt{a--z / set foundLowerCase},
\texttt{0--9 / set foundNumericDigit}, and \texttt{symbol / set foundSpecialChar}; a
transition from \textit{Scanning} to \textit{Invalid character} for anything else; and a
transition from \textit{Scanning} to \textit{Evaluate} when the input is exhausted. From
\textit{Evaluate}, one transition to \textit{Accepted} guarded by all five flags being set,
and one to \textit{Rejected} otherwise. Show \texttt{index $\geq$ 7 / set foundLongEnough}
as an action on the self-transitions.

\diagram{Password recognizer state machine}{fsm-password.png}{fsm-pass}

The five requirements, all of which must hold:

\begin{center}
\begin{tabularx}{\linewidth}{@{}p{0.26\linewidth}X@{}}
\toprule
\textbf{Requirement} & \textbf{Satisfied by} \\
\midrule
Upper case & at least one \texttt{A--Z} \\
Lower case & at least one \texttt{a--z} \\
Digit & at least one \texttt{0--9} \\
Symbol & at least one of the set listed below \\
Length & at least 8 characters \\
\bottomrule
\end{tabularx}
\end{center}

The accepted symbols are exactly these, and nothing else:

\begin{lstlisting}[basicstyle=\ttfamily\small,caption={Accepted password symbols}]
~`!@#$%^&*()_-+={}[]|\:;"'<>,.?/
\end{lstlisting}

Two rejections happen before those checks. An empty password is rejected outright, and any
character outside the four accepted classes halts the machine immediately and is reported as
an invalid character rather than as a weak password.

\textbf{A space is not an accepted character, so passphrases are rejected.} This was a
deliberate decision: the recognizer rejected spaces while password testbed case~3 asserted
that a password containing one was valid. The recognizer was confirmed to be correct, the
test case was corrected, and a new case was added to pin the rule down. Because a space is
invisible in the list of accepted symbols, \cls{PasswordCheck} reports it by name rather
than showing the generic message.

\newpage
%=======================================================================
\section{Detailed design: key flows}
\label{sec:seq}
%=======================================================================

Four flows carry the requirements that matter. Each specification below lists the lifelines
and the ordered messages, so the Astah sequence diagram is a transcription.

\subsection{Logging in}

\astah{Sequence diagram --- login}
Lifelines: \textbf{User}, \cls{ViewUserLogin}, \cls{ControllerUserLogin}, \cls{Database},
\cls{ViewAdminHome}, \cls{ViewMultipleRoleDispatch}, \cls{ViewResetPassword}.

\begin{enumerate}[leftmargin=1.4em]
  \item User enters a username and password and presses \textbf{Log in}.
  \item \cls{ViewUserLogin} $\rightarrow$ \cls{ControllerUserLogin}: \cls{doLogin(Stage)}.
  \item Controller reads both fields straight off the View.
  \item Controller $\rightarrow$ \cls{Database}: \cls{getUserAccountDetails(username)}.
        If false, show the error alert and stop. \textit{(alt fragment)}
  \item Controller $\rightarrow$ \cls{Database}: \cls{getCurrentPassword()}, compared to the
        entered password. If different, show the same error alert and stop.
        \textit{(alt fragment)}
  \item Controller builds a \cls{User} from the cached \cls{current...} fields.
  \item Controller $\rightarrow$ \cls{Database}: \cls{getCurrentOnetimePasswordFlag()}.
  \item \textit{alt} --- flag set: $\rightarrow$ \cls{ViewResetPassword}.
  \item \textit{alt} --- exactly one role: $\rightarrow$ that role's home page.
  \item \textit{alt} --- more than one role: $\rightarrow$
        \cls{ViewMultipleRoleDispatch}.
\end{enumerate}

Note that the same message is shown whether the username or the password was wrong, so the
page never reveals which usernames exist.

\diagram{Logging in}{seq-login.png}{seq-login}

\subsection{Setting up the first administrator}

\astah{Sequence diagram --- first administrator}
Lifelines: \textbf{User}, \cls{FoundationsMain}, \cls{ViewFirstAdmin},
\cls{ControllerFirstAdmin}, \cls{UserNameRecognizer}, \cls{PasswordCheck}, \cls{Database},
\cls{ViewUserUpdate}.

\begin{enumerate}[leftmargin=1.4em]
  \item \cls{FoundationsMain} $\rightarrow$ \cls{Database}: \cls{isDatabaseEmpty()} returns
        true.
  \item \cls{FoundationsMain} $\rightarrow$ \cls{ViewFirstAdmin}:
        \cls{displayFirstAdmin(Stage)}.
  \item User types; each field pushes its value to the Controller as it changes.
  \item User presses \textbf{Create admin account}.
  \item Controller $\rightarrow$ \cls{UserNameRecognizer}:
        \cls{checkForValidUserName(...)}. On an error message, stop. \textit{(alt)}
  \item Controller $\rightarrow$ \cls{PasswordCheck}: \cls{firstProblem(password)}. On a
        problem, clear both password fields, show it, stop. \textit{(alt)}
  \item Controller compares the two passwords. If they differ, clear both and stop.
        \textit{(alt)}
  \item Controller $\rightarrow$ \cls{Database}: \cls{register(user)} with
        \cls{adminRole = true}.
  \item Controller $\rightarrow$ \cls{ViewUserUpdate}: \cls{displayUserUpdate(...)}.
\end{enumerate}

\diagram{Setting up the first administrator}{seq-first-admin.png}{seq-first}

\subsection{Invitation to account}

\astah{Sequence diagram --- invitation and account creation}
Lifelines: \textbf{Admin}, \cls{ViewAdminHome}, \cls{ControllerAdminHome}, \cls{Database},
\textbf{New user}, \cls{ViewUserLogin}, \cls{ViewNewAccount}, \cls{ControllerNewAccount},
\cls{PasswordCheck}.

\begin{enumerate}[leftmargin=1.4em]
  \item Admin enters an email address, picks a role, presses \textbf{Send invitation}.
  \item Controller $\rightarrow$ \cls{Database}: \cls{emailaddressHasBeenUsed(...)}; if
        true, warn and stop. \textit{(alt)}
  \item Controller $\rightarrow$ \cls{Database}: \cls{generateInvitationCode(email, role)}.
  \item \textit{alt} --- null returned: the invitation was not stored, so report the failure
        and stop.
  \item Controller shows the code and calls \cls{ViewAdminHome.refreshStatus()}.
  \item New user enters the code on the login page and presses \textbf{Set up account}.
  \item \cls{ViewNewAccount} $\rightarrow$ \cls{Database}:
        \cls{getRoleGivenAnInvitationCode(code)}; if empty, show the invalid-code alert and
        return without displaying the page. \textit{(alt)}
  \item \cls{ViewNewAccount} $\rightarrow$ \cls{Database}:
        \cls{getEmailAddressUsingCode(code)}, shown on the page.
  \item New user chooses a username and password, presses \textbf{Create account}.
  \item Controller $\rightarrow$ \cls{PasswordCheck}: \cls{firstProblem(...)}, then the
        passwords-match check. \textit{(alt on each)}
  \item Controller $\rightarrow$ \cls{Database}: \cls{register(user)} with the role the
        invitation granted.
  \item Controller $\rightarrow$ \cls{Database}:
        \cls{removeInvitationAfterUse(theInvitationCode)}, so the code cannot be reused.
  \item Controller $\rightarrow$ \cls{ViewUserUpdate}.
\end{enumerate}

\diagram{Invitation to account creation}{seq-invitation.png}{seq-invite}

\subsection{Changing a user's roles}

\astah{Sequence diagram --- role change on Manage Users}
Lifelines: \textbf{Admin}, \cls{ViewListUsers}, \cls{ControllerListUsers}, \cls{Database},
\cls{ModelListUsers}.

\begin{enumerate}[leftmargin=1.4em]
  \item Admin selects a row. \cls{ViewListUsers} calls \cls{showSelection(user)}.
  \item View $\rightarrow$ \cls{ModelListUsers}: \cls{formatRoles(user)} for the chips.
  \item View $\rightarrow$ \cls{Database}: \cls{getNumberOfAdmins()} and
        \cls{isFoundingAdmin(...)} to decide which controls to disable.
  \item Admin ticks or unticks a role box.
  \item View $\rightarrow$ \cls{ControllerListUsers}:
        \cls{performRoleChange(username, role, value)}.
  \item Controller $\rightarrow$ \cls{Database}: \cls{updateUserRole(...)}.
  \item View calls \cls{refreshTable(username)}, which re-reads every account and reselects
        the same row.
\end{enumerate}

Draw the guard rails as an \textit{opt} fragment before step~4: the Admin checkbox and the
Delete button are disabled when the selected account is the last administrator or the
founding administrator. Add a note that the View re-reads from the database after every
change, so the checkboxes always show what was actually stored rather than what was clicked.

\diagram{Changing a user's roles}{seq-role-change.png}{seq-role}

\newpage
%=======================================================================
\section{Data model}
\label{sec:data}
%=======================================================================

All state lives in an embedded H2 database at \texttt{jdbc:h2:\textasciitilde/FoundationDatabase},
created by \cls{Database.createTables()} on first connection. There are two tables and no
foreign key between them --- an invitation is linked to the account it later creates only by
the email address, and nothing enforces that.

\astah{ER diagram or class diagram with two entities}
Two entities with the columns below. There is no declared relationship to draw; add a note
recording that \texttt{InvitationCodes.emailAddress} corresponds to
\texttt{userDB.emailAddress} by convention only, and that the row is deleted once the
invitation is used.

\diagram{Database schema}{data-er.png}{data-er}

\textbf{\texttt{userDB}} --- one row per account.

\begin{center}
\begin{tabularx}{\linewidth}{@{}llX@{}}
\toprule
\textbf{Column} & \textbf{Type} & \textbf{Notes} \\
\midrule
\texttt{id} & \texttt{INT} & Primary key, auto increment. Creation order, and therefore the
only way to identify the founding administrator \\
\texttt{userName} & \texttt{VARCHAR(255)} & Unique \\
\texttt{password} & \texttt{VARCHAR(255)} & Clear text --- issue~\ref{iss:plaintext} \\
\texttt{firstName} & \texttt{VARCHAR(255)} & \\
\texttt{middleName} & \texttt{VARCHAR(255)} & Nullable \\
\texttt{lastName} & \texttt{VARCHAR(255)} & \\
\texttt{preferredFirstName} & \texttt{VARCHAR(255)} & Nullable \\
\texttt{emailAddress} & \texttt{VARCHAR(255)} & \\
\texttt{adminRole} & \texttt{BOOL} & Default false \\
\texttt{contributorRole} & \texttt{BOOL} & Default false \\
\texttt{viewerRole} & \texttt{BOOL} & Default false \\
\texttt{curatorRole} & \texttt{BOOL} & Default false \\
\texttt{isOnetimePassword} & \texttt{BOOL} & Default false. Forces a reset at next login \\
\bottomrule
\end{tabularx}
\end{center}

\textbf{\texttt{InvitationCodes}} --- one row per outstanding invitation.

\begin{center}
\begin{tabularx}{\linewidth}{@{}llX@{}}
\toprule
\textbf{Column} & \textbf{Type} & \textbf{Notes} \\
\midrule
\texttt{code} & \texttt{VARCHAR(10)} & Primary key. Six characters from a random UUID \\
\texttt{emailAddress} & \texttt{VARCHAR(255)} & Who the invitation is for \\
\texttt{role} & \texttt{VARCHAR(20)} & Widened from 10; \texttt{Contributor} is eleven
characters and did not fit --- issue~\ref{iss:varchar} \\
\bottomrule
\end{tabularx}
\end{center}

\subsection{Identifying the founding administrator}

There is no column marking the original administrator. It is derived instead: the row with
the lowest \texttt{id}. That is sound because \cls{ViewFirstAdmin} only runs when the
database is completely empty, so the first account created is always an administrator.

The lowest \texttt{id} overall is used rather than the lowest \texttt{id} holding the admin
role, because the former never moves. If it were the latter, demoting one administrator
would silently transfer the protection to somebody else.

\begin{lstlisting}[style=sql,caption={Identifying the founding administrator}]
SELECT userName FROM userDB ORDER BY id ASC LIMIT 1
\end{lstlisting}

\newpage
%=======================================================================
\section{Traceability}
\label{sec:trace}
%=======================================================================

Each requirement below maps to the architectural element that carries it, the class that
implements it, and the automated tests that cover it. Test identifiers refer to
\cls{listUsersTesting.ListUsersTestingAutomation} unless marked \textit{PW} (the password
testbed) or \textit{UN} (the username testbed).

\begin{center}
\footnotesize
\begin{longtable}{@{}p{0.27\linewidth}p{0.17\linewidth}p{0.25\linewidth}p{0.21\linewidth}@{}}
\toprule
\textbf{Requirement} & \textbf{Architecture} & \textbf{Class / operation} &
\textbf{Tests} \\
\midrule
\endfirsthead
\toprule
\textbf{Requirement} & \textbf{Architecture} & \textbf{Class / operation} &
\textbf{Tests} \\
\midrule
\endhead
An admin can see every account in the system &
Persistence, Manage Users &
\cls{Database.getAllUserAccounts} &
A1--A6 \\
\addlinespace
Every stored attribute survives being written and read back &
Persistence &
\cls{register}, \cls{getAllUserAccounts} &
A7--A18 \\
\addlinespace
A deleted account disappears from the list &
Persistence &
\cls{Database.deleteUser} &
A19--A21 \\
\addlinespace
A user's display name prefers their preferred first name and omits an absent middle name &
Presentation &
\cls{ModelListUsers.formatFullName} &
B1--B8, C1--C3 \\
\addlinespace
A user's roles are displayed in a fixed order with no stray spacing &
Presentation &
\cls{ModelListUsers.formatRoles} &
B9--B17, C4--C6 \\
\addlinespace
An admin can grant or remove any of the four roles &
Persistence, Manage Users &
\cls{Database.updateUserRole} &
D1--D8 \\
\addlinespace
Each role is stored independently of the others &
Data model &
\cls{userDB} role columns &
D5, D10 \\
\addlinespace
An unknown role name is rejected &
Persistence &
\cls{Database.updateUserRole} &
D11 \\
\addlinespace
The system always retains at least one administrator &
Persistence &
\cls{updateUserRole}, \cls{deleteUser} &
E4 \\
\addlinespace
The founding administrator cannot be demoted or deleted &
Persistence &
\cls{Database.isFoundingAdmin} &
E1--E8 \\
\addlinespace
A password must meet all five strength requirements &
Validation &
\cls{Model.evaluatePassword} via \cls{PasswordCheck} &
PW 1--10 \\
\addlinespace
A password may not contain a space &
Validation &
\cls{Model.evaluatePassword} &
PW 11 \\
\addlinespace
A username must start with a letter, be 4--16 characters, and use only the permitted
separators &
Validation &
\cls{UserNameRecognizer.checkForValidUserName} &
UN 1--8 \\
\bottomrule
\end{longtable}
\end{center}

\subsection{Test suite summary}

\begin{center}
\begin{tabularx}{\linewidth}{@{}llX@{}}
\toprule
\textbf{Section} & \textbf{Cases} & \textbf{Covers} \\
\midrule
A & 21 & \cls{getAllUserAccounts}: list contents, full field round trip, deletion \\
B & 17 & Name and role formatting as pure functions \\
C & 6 & The formatters applied to accounts actually read from the database \\
D & 11 & \cls{updateUserRole} for each of the four roles, and rejection of unknown roles \\
E & 8 & Protection of the founding administrator \\
\midrule
\textbf{Total} & \textbf{63} & all passing \\
\bottomrule
\end{tabularx}
\end{center}

The suite runs against the real database. It never assumes the table is empty and never
asserts an absolute row count: it records a baseline, adds three uniquely named disposable
accounts, asserts on the difference, and removes them in a \cls{finally} block.

\newpage
%=======================================================================
\section{Known issues and decisions}
\label{sec:issues}
%=======================================================================

Recorded here rather than smoothed over in the diagrams, because a diagram that disagrees
with the code is worse than no diagram.

\subsection{Resolved defects}

\paragraph{Role columns that did not exist.}
\label{iss:rolecols}
\cls{updateUserRole} wrote Contributor and Viewer to a column named \texttt{newRole1} and
Curator to \texttt{newRole2}. Neither column exists in \texttt{userDB}. The resulting
\cls{SQLException} was swallowed by a \cls{catch} block that simply returned false, so
granting any non-admin role silently did nothing. Contributor and Viewer also shared the
same column name, so they would have collided even had the columns existed.
\textbf{Fixed}; covered by tests D1--D11, which fail on the original code and pass on the
corrected code.

\paragraph{Invitation codes were never consumed.}
\label{iss:invite}
\cls{ControllerNewAccount} passed \cls{ViewNewAccount.text\_Invitation.getText()} to
\cls{removeInvitationAfterUse}. That field was declared but never populated and never
displayed, so the call always received an empty string and deleted nothing. Every invitation
code therefore remained valid permanently and could create any number of accounts.
\textbf{Fixed} by using \cls{theInvitationCode}, the value that actually brought the user to
the page.

\paragraph{The role column was too narrow.}
\label{iss:varchar}
\texttt{InvitationCodes.role} was \texttt{VARCHAR(10)}, but \texttt{Contributor} is eleven
characters, so inviting a Contributor failed at the insert. The exception was caught and the
method returned the code regardless, so the administrator was shown an invitation code that
had never been stored and the recipient was told it was invalid. \textbf{Fixed}: the column
is now \texttt{VARCHAR(20)} and the method returns null on failure so the page can report
it.

\paragraph{Stale page data.}
\label{iss:stale}
\cls{ViewAdminHome} set the signed-in username and both status counts in its constructor.
Because the class is a singleton, the constructor runs once, so after the first
administrator signed in the page kept showing that person's name and the counts as they
were at that moment. \textbf{Fixed} by moving the refresh into \cls{displayAdminHome}.

\paragraph{The founding administrator could be removed.}
\label{iss:founder}
The last-administrator rule only prevented the system from ending up with no administrators.
Once a second administrator existed, that second account could demote or delete the
original one. \textbf{Fixed} by protecting the founding account in \cls{Database}, so both
the Manage Users page and the older Add/Remove Roles page inherit the guard.

\subsection{Open issues}

\paragraph{Passwords are stored in clear text.}
\label{iss:plaintext}
\texttt{userDB.password} holds the password as typed, and login compares strings directly.
Anyone who can read the database file can read every password. The account page previously
displayed it on screen; it is now masked, but the storage is unchanged. Fixing this properly
means salted hashing and is a change to \cls{register}, \cls{updatePassword} and every login
method.

\paragraph{The schema is never migrated.}
\label{iss:migrate}
\cls{createTables} uses \cls{CREATE TABLE IF NOT EXISTS}, which does nothing at all when the
table already exists. A database created before the roles changed keeps the old columns, and
account creation then fails outright with \textit{Column ``CONTRIBUTORROLE'' not found}. The
same applies to the widened \texttt{role} column. \textbf{Anyone whose database file predates
those changes must delete} \texttt{FoundationDatabase.mv.db} \textbf{before running the
application.}

\paragraph{Exceptions are swallowed.}
Several \cls{Database} methods catch \cls{SQLException} and return a default rather than
reporting the failure. Two of the defects above were invisible for exactly this reason: the
database refused the operation and the caller could not tell the difference between refusal
and success.

\paragraph{A username cannot be changed.}
The account page shows the username as read-only. The original page offered an
\textit{Update Username} button wired to nothing, which invited an action that could never
work; the button was removed rather than left as a dead control.

\subsection{Decisions recorded}

\begin{center}
\begin{tabularx}{\linewidth}{@{}p{0.30\linewidth}X@{}}
\toprule
\textbf{Decision} & \textbf{Reasoning} \\
\midrule
Passwords may not contain spaces &
The recognizer rejected spaces while a testbed case asserted a password containing one was
valid. The recognizer was confirmed correct and the test case corrected. \\
\addlinespace
The founding administrator is identified by lowest \texttt{id} &
That value never moves, whereas ``lowest id holding the admin role'' would transfer the
protection to someone else as soon as an administrator was demoted. \\
\addlinespace
Role and delete actions were added to Manage Users &
Acting on the row already on screen removes the class of error where the wrong name is
picked from a dropdown. The older single-purpose pages were kept so nothing that depended on
them broke. \\
\addlinespace
Formatting logic lives in \cls{ModelListUsers} &
The page and the automated tests call the same code, so the two cannot disagree. Testing a
copy would let the page drift without any test noticing. \\
\bottomrule
\end{tabularx}
\end{center}

\end{document}
